\section{Discussion}

The proposed PreForceVision framework defines a stepwise route for evaluating whether preload-induced bolt-head deformation can be mechanically observed, optically extracted, and related to preload with sufficient repeatability.
%The technical plausibility of the framework depends on the signal-to-noise ratio of the deformation measurement. The preliminary FEM, strain-gauge, and optical 3D results indicate that preload produces a measurable bolt-head deformation in the order of 10 µm to 20 µm. However, this deformation magnitude is also the main limitation. It cannot be assumed that ordinary images under uncontrolled field conditions contain sufficient information. The proposed method therefore requires controlled acquisition, reference-loaded image pairs, careful registration, suitable surface texture, and uncertainty-based rejection of unreliable measurements.

\subsection{Advantages}
The main advantage is mechanical relevance. The target quantity is preload, not only visible looseness or rotation. This links visual inspection to the functional variable governing joint functional behavior. The second advantage is multimodal validation. Finite-element analysis, strain gauges, optical 3D measurement, and DIC provide complementary evidence and reduce the risk of training a model on non-physical image artifacts. The third advantage is interpretability. Deformation descriptors such as radial displacement, radial strain, center-to-edge contrast, and symmetry can be inspected before and after regression. This is important because the deformation signal is small and requires physically traceable evaluation. 
%This is preferable to a purely black-box image classifier, especially when the optical signal is small.

\subsection{Limitations}
The main limitation is the small signal magnitude. The preliminary deformation is in the micrometer range. Although this magnitude is measurable with high-resolution optical 3D systems, recovering it from image pairs requires controlled lighting, stable camera positioning, suitable surface texture, accurate registration, and sufficient image resolution. A second limitation is surface condition. Speckle patterns or matte coatings improve deformation extraction but reduce practical convenience, while untreated metallic heads are more affected by reflection and weak texture. A third limitation is reference dependence. The first-stage method requires a reference image of the same bolt-head surface; absolute preload estimation from a single image would require stronger assumptions and a substantially larger dataset. A fourth limitation is generalization. A model trained on one bolt type, surface condition, and acquisition setup cannot be assumed to generalize without explicit validation.
%The main limitation is signal magnitude. The preliminary deformation is on the order of tens of micrometers. Such deformation can be measured by optical 3D systems, but its recovery from ordinary images requires controlled texture, lighting, registration, and focus. A second limitation is surface condition. Speckle patterns and matte coatings improve DIC, but they reduce practical convenience. Untreated metallic heads are more realistic but more vulnerable to specular reflection and weak texture. A third limitation is reference dependence. The proposed first-stage method requires a reference image; absolute preload estimation from a single image would require much stronger assumptions and a substantially larger data set. A fourth limitation is generalization. A model trained on a single bolt type, surface condition, and camera configuration may not generalize without explicit validation.

\subsection{Implications}

The Preliminary study is a mechanically validated optical feasibility study. The recommended sequence is: establish the deformation field using FEM and strain gauges; confirm the central depression using optical 3D measurements; extract DIC or optical-flow features from controlled image pairs; verify monotonicity and repeatability; then train and evaluate regression models. This sequence establishes the scientific basis before evaluating low-cost or field-compatible acquisition.